% TOP_PROBLEM: {"id": 30007504, "title": "Finite assembly of the height-three chromatic overlap", "category_id": 7, "category": {"id": 7, "name": "topology", "display_name": "Topology", "description": "Properties preserved under continuous deformations.", "slug": "topology", "order_index": 7, "created_at": "2026-07-31T15:26:25.671Z"}, "status": "open"}
% FIRST_POSTED: 2026-10-06
% ENTRY_KIND: statement_only
% !TeX program = lualatex
\documentclass[11pt]{article}
\usepackage[margin=25mm]{geometry}
\usepackage{fontspec}
\setmainfont{Latin Modern Roman}
\usepackage{amsmath,amssymb,mathtools}
\usepackage{unicode-math}
\setmathfont{Latin Modern Math}
\usepackage{xurl,hyperref}
\hypersetup{colorlinks=true,urlcolor=blue}
\setcounter{secnumdepth}{0}
\setlength{\parindent}{0pt}
\setlength{\parskip}{5pt}
\setlength{\emergencystretch}{3em}
\newcommand{\R}{\mathbb R}
\newcommand{\N}{\mathbb N}
\newcommand{\Z}{\mathbb Z}
\newcommand{\Q}{\mathbb Q}
\newcommand{\C}{\mathbb C}
\newcommand{\D}{\mathbb D}
\newcommand{\F}{\mathbb F}
\newcommand{\sref}[2]{\hyperlink{source:#1:#2}{[S#2]}}
\newcommand{\cataloguescope}[1]{#1}
\begin{document}
\section*{Finite assembly of the height-three chromatic overlap}\label{30007504}
\subsection{Definitions and mathematical statement}
\paragraph{Statement recorded in the supplied source.}
At $p=3$, can $L_{2}L_{K(3)}S_{3}^\wedge $ be built from $L_{0}S_{3}^\wedge $, $L_{1}S_{3}^\wedge $ and $L_{2}S_{3}^\wedge $ by finitely many shifts, sums, cofibers and retracts in the category of E(2)-local modules over the derived three-complete sphere?
\subsection{Short English statement}
At $p=3$, can $L_{2}L_{K(3)}S_{3}^\wedge $ be built from $L_{0}S_{3}^\wedge $, $L_{1}S_{3}^\wedge $ and $L_{2}S_{3}^\wedge $ by finitely many shifts, sums, cofibers and retracts in the category of E(2)-local modules over the derived three-complete sphere?
\subsection{Sources}
\begin{itemize}
\item[S1] Original question source or mathematical context.
\url{https://arxiv.org/abs/2402.00960}.
\item[S2] Original question source or mathematical context.
\url{https://doi.org/10.1007/s00222-019-00943-9}.
\item[S3] Original question source or mathematical context.
\url{https://doi.org/10.4171/owr/2015/14}.
\item[S4] Original question source or mathematical context.
\url{https://www-users.cse.umn.edu/~tlawson/hovey/}.
\item[S5] OpenAI, manuscript citation page. Source of the formulation motivating this local problem.
\url{https://github.com/openai/math/blob/adc7f1241b42e322a6451854ab7e4b4c146bf78a/preprints/Failure-of-finite-assembly-for-a-chromatic-overlap-at-the-prime-three-September-25-2026/README.md}.
\end{itemize}
\end{document}
